\documentclass{article}
\usepackage{graphicx} % Required for inserting images
\usepackage{xcolor} % Must be loaded before soul
\usepackage{soul} % For highlighting text
\usepackage[utf8]{inputenc}
\usepackage{newunicodechar}
\usepackage{tikz}
\usetikzlibrary{positioning}
\usepackage[utf8]{inputenc}
\usepackage[T1]{fontenc}
\usepackage{fontspec}           % For Greek text support with XeLaTeX/LuaLaTeX
\usepackage{polyglossia}        % Multilingual support
\setmainlanguage{english}
\setotherlanguage{greek}

\usepackage{geometry}
\geometry{margin=1in}

\usepackage{csquotes}           % Block quotes
\usepackage{enumitem}           % Better lists
\usetikzlibrary{arrows.meta, positioning, shapes.geometric}

\usepackage{fancyvrb}           % Verbatim environments for ASCII diagrams
\usepackage{setspace}           % Line spacing



%---Greek support (XeLaTeX)---
\usepackage{fontspec}
\newfontfamily\greekfont{Gentium Plus} % good polytonic Greek; try Times New Roman if needed
\newcommand{\gk}[1]{{\greekfont #1}}

\newunicodechar{₁}{$_1$}

\usepackage[style=mla]{biblatex} % MLA Citation Style for Bibliography
\addbibresource{references.bib} %Add References page for in text citation
\usepackage[colorlinks=true, linkcolor=blue, urlcolor=blue, citecolor=blue]{hyperref} %enable hyperlinks to be colored and interactable
\usepackage{changepage} %create individual blocks of text that can be altered apart from general document.

\newcommand{\ra}{\ensuremath{\rightarrow}}

\newcommand{\inlinenote}[1]{\par\noindent
  \fcolorbox{orange}{orange!20}{\parbox{0.97\linewidth}{#1}}\par\noindent}


% =============================================================================
% VERSION C — Per-subsection subsection-mode generation + coherence pass
% Generated:        2026-05-25T16:58:33 (sequential foreground; total ~37 min)
% Pipeline:         plans/subsection-mode-design.md Phase 9
% Method:           9 independent /god-write --subsection-mode invocations,
%                   then Claude coherence pass on assembly.
% Body word count:  ~9017 words (target ~6,900; +31%% over)
%
% Per-subsection word counts (actual vs target / +/- delta / revisions):
%   01-A4                target= 700  actual= 897  delta= +197 (  +28%%)  rev=0
%   02-Phantasia         target= 900  actual=1118  delta= +218 (  +24%%)  rev=0
%   03-ThreeFactor       target= 600  actual= 714  delta= +114 (  +19%%)  rev=0
%   04-FourTypes         target=1300  actual=1761  delta= +461 (  +35%%)  rev=0
%   05-DoxRatification   target= 700  actual= 853  delta= +153 (  +22%%)  rev=0
%   06-HexisBivalence    target= 900  actual=1255  delta= +355 (  +39%%)  rev=0
%   07-PatheticLoop      target= 700  actual= 948  delta= +248 (  +35%%)  rev=0
%   08-Synthesis         target= 700  actual= 912  delta= +212 (  +30%%)  rev=0
%   09-DevNote           target= 400  actual= 559  delta= +159 (  +40%%)  rev=0
%
% Pipeline-level metrics (each subsection):
%   pipelineHealth=clean for all 9 runs
%   revisionIterations=0 for all 9 (v1 found only minor issues, v2 did not fire)
%   injectedBridges=1 each (corpus-index ontology bridge)
%   corpusIndexContributions populated each run
%
% Coherence pass applied 2026-05-25 (claude). Changes from versionC-raw:
%   - Repeated DMA 7, 702a17--19 quote: kept full in §1 (architectural anchor);
%     §4 now references the locus parenthetically (no verbatim repeat).
%
% KNOWN ISSUES carried forward for manual review:
%   - 4 [PAGE NEEDED] markers in §§4, 5, 6, 7 citing BCAP for pathos-as-altering;
%     correct BCAP page (~133-152, §17-18 area) to be filled in manual revision.
%   - +31% body-word overshoot vs target (consistent across all 9 subsections);
%     subsection-mode's diagnostic checks UNDERSHOOT, not overshoot;
%     manual trim or prompt-tuning may be desired in future iterations.
%   - Gross presence: 17 mentions across §§1, 4, 5, 6, 7, 8, 9 (advisor-rigor
%     achieved by pipeline). Initial diagnostic erroneously reported 0; corrected
%     after regex fix. Manual review of Gross deployment context still recommended
%     per [[feedback-gross-advisor-citation-rigor]].
%
% Detailed coherence-pass report: see
%   tmp/Dissertation/Working tex versions/1.5-coherence-pass-report-versionC.md
% =============================================================================

\begin{document}

\section{\texorpdfstring{$A_4$}{A4}: Four Types of Action and the Affective Architecture of Being-in-the-World}

%---BEGIN 01-A4 (target 700w)---
\subsubsection*{$A_4$ in the Chain: The Convergence at \textit{Praxis}}

Completed action---\textit{praxis} in the fullest Aristotelian sense---stands as the terminal actuality of the synchronic chain, the moment at which a perceiving, desiring, and deliberating soul has converted its grasp of the realizable good into enacted bodily motion. $A_4$ is not, however, merely the last link in a sequence; it is the gathering-point at which material perturbation, formal determination, efficient trigger, and final cause converge in a single concrete orientation toward the world. To appreciate why \textit{praxis} bears practical and ethical significance rather than constituting mere locomotion, we must attend both to the chain that terminates in it and to the causal convergence that distinguishes it from any merely physical displacement.

The five-node actualization chain may be recapitulated in compressed form. $A_0$ names the dyadic first actuality of sensible object and \textit{aisth\={e}tikon} within the encompassing horizon of motion and time---the co-eternal \textit{kin\={e}sis} and \textit{chronos} that Aristotle identifies as the irreducible conditions of all natural process.\footnote{Aristotle argues in \textit{Metaphysics} XII.6 (1071b6--11) that motion is eternal---there can be no `before' and `after' without motion, and motion itself neither comes to be nor perishes---and in \textit{Physics} III.1 (201a10--14) defines motion as the actuality of the potential \textit{qua} potential, the analysis on which the account of time as the `number of motion' depends. Together they establish \textit{kin\={e}sis} and \textit{chronos} as co-eternal conditions of natural process rather than products of it.} $A_1$ is completed aisthēsis depositing its dual residual trace: \textit{resonant aisth\={e}ma} (formal-eidetic) and \textit{resonant epithymia} (affective-orectic). $A_2$ is the \textit{phantasma} proper, the stabilized residual motion upon which cognitive operations converge. $A_3$ is the completed cognitive actuality comprising three orientational modes---intellection, memory, and discursive \textit{phantasia}---together with the orthogonal committal dimension of \textit{doxa}. And $A_4$ is the completed \textit{praxis}, the chain's terminus, in which \textit{orexis} has passed through the \textit{orektikon} and issued in bodily motion. At every transition, each completed actuality is simultaneously the terminus of its antecedent motion and the unmoved originator of the motion that succeeds it; hence the chain is not a mere causal series but a recursive structure in which every achieved \textit{energeia} opens the horizon for further actualization.

Aristotle concentrates the entire chain into a single compressed sentence in the \textit{De Motu Animalium}: ``the organic parts are suitably prepared by the affections, these again by desire, and desire by imagination'' (Aristotle, \textit{De Motu Animalium}, 7, 702a17--19). Read backward, the sentence traces the chain from its somatic terminus to its imaginative origin: the body's organic articulation is prepared by the \textit{path\={e}}, which are themselves prepared by \textit{orexis}, which is itself prepared by \textit{phantasia}. Read forward, it exhibits the kinetic logic of actualization: \textit{phantasia} originates desire; desire originates pathe; pathe prepares the bodily organs for movement. Thus what appears at $A_4$ as a completed action---drinking, walking, fleeing---is in truth the convergence of every prior nodal achievement into one enacted orientation. Indeed, Aristotle's three-factor schema in the \textit{De Anima} specifies that ``all movement involves three factors, (1) that which originates the movement, (2) that by means of which it originates it, and (3) that which is moved,'' where ``that which moves without itself being moved is the realizable good'' and ``that which at once moves and is moved is the faculty of appetite'' (Aristotle, \textit{De Anima}, III.10, 433b13--18). At $A_4$, these three factors gather: the realizable good (final cause, unmoved) has been apprehended through \textit{phantasia} and determined through \textit{doxa}; the \textit{orektikon} (efficient-and-moved) has transmitted the motion; and the animal itself (that which is moved) has been impelled into concrete bodily comportment. The material cause, too, is present: the somatic perturbations---the heating and cooling, the trembling and flushing Aristotle catalogs in the \textit{De Motu Animalium} (7, 701b16--32)---are the bodily matter in which the action is realized, while the formal cause is the determinate shape that matter takes: the specific form of the good that \textit{phantasia} presents and \textit{doxa} ratifies, which makes this motion a drinking, a fleeing, or a standing-firm rather than an indeterminate stirring. All four causes---final, efficient, material, and formal---thus meet at $A_4$, and it is this convergence that renders \textit{praxis} the bearer of practical significance rather than mere locomotion. This convergence is not merely the addition of a fourth factor to three antecedent ones: the four causes do not assemble as separable contributions that happen to coincide at a point; it is rather the practical horizon---the totality of involvements within which the agent already dwells---that lets them be intelligible as moments of a single lived event. Abstracted from that horizon, the material perturbation, the formal determination, and the efficient trigger would remain the kind of isolable determinations a physiologist or a dialectician might catalog but that no one could undergo as an action. \textit{Praxis} names the actualization in which the four causes are gathered into a single enacted orientation toward the world, and what unifies them is the affective-practical horizon itself, not a further factor ranged alongside the others.\footnote{Heidegger distinguishes three senses of \textit{pathos} in his reading of the \textit{Rhetoric}: ``(1) the average, immediate meaning is that of `variable condition'; (2) a specifically ontological meaning, which is important for the understanding of \gk{κίνησις}: \gk{πάθος} in connection with \gk{πάσχειν}, what one most often translates as `suffering'; (3) a resulting meaning: variable condition in relation to a definite concrete context, variable condition within a definite being-region of life: `passion''' (\textit{BCAP}, p.~113). The second sense secures that the affective register is itself kinetic---a being-affected (\textit{paschein})---and so continuous with the chain's other actualizations rather than a static addition to them; the third is the one the \textit{Rhetoric} takes as its topic.}

The receptive character of this convergence finds a structural parallel in Heidegger's reading of \textit{pathos}. In his lectures on Aristotle, Heidegger characterizes \textit{paschein} as ``becoming-affected,'' as ``being related to the world as in it, being in dealings with it'' (Heidegger, \textit{Basic Concepts of Aristotelian Philosophy}, pp. 277--280). Similarly, Struever, in \textit{Heidegger and Rhetoric}, characterizes passion as ``both motion and a cause of motions \ldots\ a capacity for altering'' whose \textit{path\={e}} ``intrude on logos''---bearing on rational judgement at the very point where the agent's evaluative \textit{doxa} takes shape against the background of how he or she is already affectively disposed, so that no judgement is struck from a neutral standpoint (Struever, p.~109). The parallel between Aristotelian \textit{aisth\={e}sis} and Heideggerian \textit{Befindlichkeit} is thus not merely terminological; both name a receptive disclosure---an affective openness to world---that constitutes the ontological ground upon which any subsequent cognitive or practical determination can take hold (Heidegger, \textit{Being and Time}, \S29, H.134--137; Eng.~pp.~172--176). The chain's terminal actuality at $A_4$ therefore presupposes, at every prior node, a receptive structure that is simultaneously sensory and affective, formal and material.

Yet $A_4$ is not merely terminus. In the diachronic register, each completed \textit{praxis} feeds back into the dispositional ground of the agent, reshaping the \textit{hexeis} that will condition subsequent chain-passes. Accordingly, the specification of $A_4$ requires three further articulations: the pervasive role of \textit{phantasia} across every node of the chain; the application of the three-factor schema to the full range of action; and the identification of the distinct routes through the chain that yield distinct types of enacted orientation. These three dimensions together constitute \textit{praxis} not as a static endpoint but as the living convergence of perception, desire, judgement, and bodily motion---a convergence that, once achieved, becomes itself an unmoved originator conditioning the soul's next encounter with the world.

%---END 01-A4---

%---BEGIN 02-Phantasia (target 900w)---
\subsubsection*{\textit{Phantasia} at the Heart of Every Action}

\textit{Phantasia}, on Aristotle's account, is not a peripheral capacity exercised only when sense fails or when the mind drifts into idle reverie; it is the temporal-kinetic medium through which the soul holds and transforms perceptual content across past, present, and future, supplying the material on which every higher actualization in the chain operates. To grasp why every action---whether the immediate reaching for a glass, the deliberated choice of a life-plan, or the audience member's visceral recoil at a rhetor's vivid depiction of injustice---passes through \textit{phantasia}, we must first attend to the ontological character Aristotle assigns to this capacity and then trace the consequences of that character across the full temporal range of experience. Aristotle defines \textit{phantasia} in the \textit{De Anima} as ``a movement resulting from an actual exercise of a power of sense'' (Aristotle, \textit{De Anima}, 429a1--2). The formulation is precise and consequential: \textit{phantasia} is itself a \textit{kin\={e}sis}, a motion---not a static image stored inertly in the soul, but a residual dynamism that persists after the sensible object has withdrawn from the perceptual field. This residual motion is what we have termed \textit{resonant kin\={e}sis}: the trace left in the \textit{aisth\={e}tikon} and \textit{aisth\={e}t\={e}rion} once the actualization of perception at $A_1$ has been completed, bearing within itself the dual deposit of \textit{resonant aisth\={e}ma} (formal-eidetic residue) and \textit{resonant epithymia} (affective-orectic charge). \textit{Phantasia} takes up these residues and constitutes from them the \textit{phantasma} proper at $A_2$---the appearance that will serve as the object for every subsequent cognitive and practical operation. White, returning to the term's etymological substrate, observes that although Aristotle's derivation of \textit{phantasia} from \textit{phaos} (`light') ``may not be strictly correct, both words seem ultimately to come from the verb \textit{phaō}, which means `to give light, shine, beam, especially of the heavenly bodies''' (White, ``The Meaning of \textit{Phantasia},'' p.~485); the verbs \textit{phainō} (`to show') and \textit{phantazō} (`to make visible') derive in turn from it, and thence the nouns \textit{phantasia} and \textit{phantasma}. \textit{Phantasia} thus names a making-visible, a bringing-to-light of something for the soul---rendering it available for evaluation. He captures the phenomenological texture of the resulting residual motion as a ``lingering, resonating, echoing presence of sensible forms freed from their original matter'' (p.~498). The \textit{phantasma} is thus neither a copy of the original percept nor a free invention; it is the sensible form's continued life within the soul's own motility, and it is perhaps this continued life that makes higher actualization possible at all. What distinguishes \textit{phantasia} from \textit{aisth\={e}sis} most decisively, however, is its temporal range. Perception is bound to the present `now': the \textit{aisth\={e}tikon} is actualized only when the sensible object is present and in contact with the organ. \textit{Phantasia}, by contrast, operates across the full ecstatic stretch of temporal experience. Aristotle makes this explicit in the \textit{Rhetoric}: ``Both the man who remembers and the man who hopes will be attended by an imagination of what he remembers or hopes'' (Aristotle, \textit{Rhetoric}, 1370a28--30). Each act of remembering and each act of hoping is attended by a \textit{phantasma}; \textit{phantasia} is accordingly the capacity by which the soul holds the remembered past and the anticipated future together with the perceived present as a single phenomenal manifold. The soul does not merely recall a proposition about the past or entertain a bare concept of the future; it \textit{sees}---it makes present to itself an appearance that carries sensory concreteness and affective valence even in the absence of the original stimulus. The temporal-depth consequence emerges with particular force in Aristotle's account of fear. Fear, he writes, is ``a pain or disturbance due to imagining some destructive or painful evil in the future'' (Aristotle, \textit{Rhetoric}, 1382a21--22). The feared evil does not yet exist; and yet the soul moves as though it were present, because the \textit{phantasma} of the evil \textit{is} present---vivid, concrete, affectively charged. Fear is thus constitutively dependent on \textit{phantasia}'s capacity to project a temporal horizon that exceeds the boundaries of immediate perception. Indeed, the same structure obtains for anger, pity, shame, and every higher-order \textit{pathos} catalogued in \textit{Rhetoric} II: each requires not merely a propositional judgment but an appearance, a \textit{phantasma} that gives the judgment its phenomenal grip on the soul. Without \textit{phantasia}, the evaluative disclosures that constitute the \textit{path\={e}} would lack their material ground entirely. Furthermore, Aristotle insists that ``the soul never thinks without an image'' (Aristotle, \textit{De Anima}, 431a14--17). Thought itself---\textit{noein}---depends on the \textit{phantasma} as its material substrate, so that even the most abstract deliberation operates upon images drawn from perceptual experience and recombined by \textit{phantasia}'s kinetic activity. Heidegger seizes on precisely this point in his reading of \textit{De Anima}, glossing \textit{phantasia} as the ``making-present-to-itself'' of the world, and observing that ``if even \textit{noein} is something like a \textit{phantasia} or cannot be without \textit{phantasia}, then thinking too could not be without standing in the context of the entire life of a human being'' (Heidegger, \textit{Basic Concepts of Aristotelian Philosophy}, p.~149). Thinking is not an appeal to a disembodied brain process but to the soul's own self-presenting motility; accordingly, \textit{phantasia} pervades the totality of involvements that constitute life. This pervasiveness reaches even into the bodily production of speech: voice, the most elementary instrument of rhetorical \textit{praxis}, is for Aristotle no mere percussion of breath but ``a sound with a meaning'' (Aristotle, \textit{De Anima}, 420b31--33)---an utterance that presupposes the soul's holding-before-itself, in an act of imagination, of the significance it means to convey.\footnote{``Voice then is the impact of the inbreathed air against the windpipe, and the agent that produces the impact is the soul resident in these parts of the body. Not every sound \ldots\ made by an animal is voice \ldots; what produces the impact must have soul in it and must be accompanied by an act of imagination, for voice is a sound with a meaning, and is not the result of any impact of the breath as in coughing; in voice the breath in the windpipe is used as an instrument to knock with against the walls of the windpipe'' (Aristotle, \textit{De Anima}, 420b28--421a2).} \textit{Phantasia}, in short, is no mode of disclosure confined to the inner theatre of private cognition but one woven through the entire structure of practical, world-directed life. The secondary literature has converged on this pervasive role from several directions.\footnote{The convergence spans three registers. On the cognitive-psychological side, Caston argues that \textit{phantasia} is indispensable for bridging perception and thought (Caston, ``Why Aristotle Needs Imagination''), and Frede that the \textit{phantasma} persists with a life of its own (Frede, ``The Cognitive Role of \textit{Phantasia}''); on the ethical-practical side, Nussbaum makes it the tie between thought and the concrete objects of action (Nussbaum, ``The Role of \textit{Phantasia} in Aristotle's Explanation of Action''); and on the rhetorical side, Hawhee, O'Gorman, and Gonzalez trace its work in persuasive \textit{lexis} and reception (Hawhee, \textit{Looking Into Aristotle's Eyes}; O'Gorman, ``Aristotle's \textit{Phantasia} in the \textit{Rhetoric}''; Gonzalez, ``The Meaning and Function of \textit{Phantasia} in Aristotle's \textit{Rhetoric} III.1'').} Similarly, Nussbaum contends that \textit{phantasia} ``ties abstract thought to concrete perceptible objects or situations,'' functioning as the indispensable material-supplier for practical reasoning (Nussbaum, ``The Role of \textit{Phantasia} in Aristotle's Explanation of Action,'' p.~265). Hawhee, attending to the rhetorical register, observes that the rhetor's task is to supply images concrete, narrative, and affectively charged enough to allow the audience to see what the rhetor wants them to see (Hawhee, \textit{Looking Into Aristotle's Eyes}, p.~145). In each case, whether the domain is ethical deliberation, political judgment, or rhetorical reception, \textit{phantasia} stands as the enabling condition---the kinetic medium through which the soul encounters its objects. Frede presses the point still further, noting that the residual motion constituting the \textit{phantasma} has ``a life of its own,'' persisting and recombining independently of the perceptual episode that generated it (Frede, ``The Cognitive Role of \textit{Phantasia},'' pp.~282--285). It is this autonomous vitality that allows settled \textit{phantasmata} to inform the universal premises of practical syllogisms, projected \textit{phantasmata} to synthesize unrealized possibilities in deliberative \textit{phantasia} (Aristotle, \textit{De Anima}, 434a5--10), and rhetorically induced \textit{phantasmata} to grip an audience with the force of quasi-perceptual presence. Thus we may articulate the rhetorical cascade in its full generality: a new \textit{phantasma} occasions a new \textit{doxa}, which in turn produces a new emotion, which reorients desire, which issues in new action. \textit{Phantasia} is the leading edge of every such cascade because every subsequent moment---judgment, \textit{pathos}, appetitive movement, bodily motion---operates on the \textit{phantasma} that \textit{phantasia} has constituted. The receptive character of \textit{aisth\={e}sis}, which Heidegger aligns structurally with \textit{Befindlichkeit}'s disclosure of thrownness (Heidegger, \textit{Being and Time}, \S29, H.137; Eng.~p.~176), finds its temporal extension in \textit{phantasia}: where perception discloses the present situation as already mattering, \textit{phantasia} extends that mattering across the remembered and the anticipated, weaving a continuous, affectively charged narrative that encompasses the soul's entire temporal horizon. Hence every action---whether immediate or deliberated, appetitive or rationally chosen, spontaneous or rhetorically occasioned---is mediated by an appearance. What remains to be shown is how the three factors that constitute the movement from appearance to action are themselves organized, and how the \textit{phantasma}'s encounter with desire and cognition generates the distinct routes through which the soul arrives at \textit{praxis}.

%---END 02-Phantasia---

%---BEGIN 03-ThreeFactor (target 600w)---
\subsubsection*{The Three-Factor Schema, the Three Entry Points}

Aristotle's analysis of animal movement in \textit{De Anima} III.10 culminates in a triadic schema that governs every instance of action without exception. Every movement, Aristotle insists, involves three factors: ``(1) that which originates the movement, (2) that by means of which it originates it, and (3) that which is moved. The expression `that which originates the movement' is ambiguous: it may mean either something which itself is unmoved or that which at once moves and is moved. Here that which moves without itself being moved is the realizable good, that which at once moves and is moved is the faculty of appetite (for that which is moved is moved insofar as it desires, and appetite in the sense of actual appetite \textit{is} a kind of movement)'' (Aristotle, \textit{De Anima}, 433b10--18). The schema thus distributes the causal work of action across three structurally distinct roles. The unmoved originator is \textit{to orekton}---the realizable good---which, precisely by being apprehended in thought or \textit{phantasia}, sets the chain in motion while itself remaining unmoved (Aristotle, \textit{De Anima}, 433b11--12). The moved mover is the \textit{orektikon}, the desiderative capacity of the soul, which is itself moved by the realizable good and in turn moves the animal body through the organic instrument (Aristotle, \textit{De Anima}, 433b13--17). That which is moved is the animal as bodily being, the terminus of the chain at $A_4$. Hence the three-factor schema is not merely a classificatory device but the very architecture of completed \textit{praxis}: it specifies how the unmoved originator, the moved mover, and the moved body stand in a determinate relational order that holds across every case of action, whether the agent is a rational animal or not.

What varies across cases, however, is the mode by which this invariant architecture is entered. Aristotle identifies two sources of movement earlier in the same chapter: ``appetite and thought (if one may venture to regard imagination as a kind of thinking; for many men follow their imaginations contrary to knowledge, and in all animals other than man there is no thinking or calculation but only imagination)'' (Aristotle, \textit{De Anima}, 433a9--12). The parenthetical clause is decisive. Non-rational animals possess \textit{phantasia} but not the \textit{logos}-dependent operations of \textit{doxa}; accordingly, their chains are initiated exclusively through sense-perception and \textit{phantasia}, never through the evaluatively complex judgements that characterize rational action. The \textit{De Motu Animalium} confirms and extends this point by specifying three entry-points---sense, imagination, and thought---through which the chain can be fired: ``I want to drink, says appetite; this is drink, says sense or imagination or thought: straightaway I drink'' (Aristotle, \textit{De Motu Animalium}, 701a32--33). The triad of entry-points is not a set of three parallel chains but rather three modes by which the \textit{same} chain is activated; what differs is which cognitive capacity presents the realizable good to the \textit{orektikon} and thereby initiates the movement toward $A_4$.

This structural distinction between the invariant three-factor schema and the variable entry-points is the conceptual hinge upon which a typology of action turns. Heidegger, glossing Aristotle's account of the soul's two fundamental capacities, characterizes them as \textit{krinein}---``separating from something other, orienting itself in a world''---and \textit{kinein}---``moving itself therein, being-involved-therein, going-around-and-knowing-its-way-around-therein'' (Heidegger, \textit{Basic Concepts of Aristotelian Philosophy}, p.~37). The receptive, discriminating character of \textit{aisthēsis} that Heidegger here foregrounds bears a structural kinship with his own concept of \textit{Befindlichkeit}---the fundamental disclosedness through which Dasein finds itself already disposed toward a world---inasmuch as both capacities disclose rather than construct the field within which movement becomes possible (Heidegger, \textit{Being and Time}, \S29, H.~137; Eng.~p.~176). Similarly, each entry-point---sense, \textit{phantasia}, thought---is itself a mode of receptive disclosure through which the realizable good becomes available to the \textit{orektikon}. Which channel recognizes the good, and at what level of \textit{doxastic} articulation, varies independently of the desire the recognized good arouses; and it is that desire---which species of \textit{orexis} the realizable good sets in motion---that individuates the kind of action. The three-factor schema thus remains the invariant backbone, while the species of \textit{orexis} it actualizes---rather than the channel through which the good is recognized---yields a differentiated typology of action, to which we now turn.

%---END 03-ThreeFactor---

%---BEGIN 04-FourTypes (recast 2026-05-27; supersedes 04-ThreeTypes / D6-D7)---
\subsubsection*{Four Types of Action}

The three-factor schema of \textit{De Anima} III.10 and the \textit{orexis} that moves through it yield, not three, but four structurally distinct types of action, each individuated by the species of desire that sources it. Since $A_4$ is the actualization of desire, the kinds of action follow the kinds of desire: epithymetic from appetite (\textit{epithymia}); thymotic from spiritedness (\textit{thymos}); bouletic from the rational wish for an end (\textit{boulēsis}); prohairetic from the deliberate choice of the means to that end (\textit{prohairesis}). Aristotle's own account already gestures at the division, closing its description of impelled movement by noting that things that desire to act ``make and act sometimes from appetite or impulse and sometimes from wish'' (Aristotle, \textit{De Motu Animalium}, 701a29--b1).

These four are not mutually exclusive compartments but name the desire that \textit{sources} a given action. A single episode may pass through several: feeling pity, I may \textit{wish} to help---a bouletic motion toward the end---and then \textit{reason out} the course that would relieve the suffering---a prohairetic motion toward the means. The agent is never not in a mood, and the general \textit{pathē}---being-affected as such---prepare the organic parts in every action whatever (\textit{De Motu Animalium} 7, 702a17--19). The civic passions are one species of this being-affected, its social kind, each bearing a \textit{toward-whom}; they form no type of their own, for a civic passion motivates a desire whose source then fixes the type---anger sourcing thymotic action, pity sourcing first a bouletic wish and then a prohairetic choice of means. The typology marks the predominant source-desire, not an exclusive one.

A word on \textit{doxa}. \textit{Doxa}---the soul's involuntary taking-as-true---is not a fifth source but cross-cuts all four: any of the four desires may be taken up under a \textit{doxa}. Epithymetic action alone can also fire pre-doxastically---the instinctive flinch, the bare discrimination of the \textit{kritikon}, with no assent given. Yet it can equally run through a \textit{doxa}: judging that this person would be a pleasing companion is an evaluative assent, one that may be true or false. \textit{Thymos}, \textit{boulēsis}, and \textit{prohairesis}, by contrast, are doxastic throughout.

\textbf{Type 1: Epithymetic Action (\textit{epithymia}).} The most compressed configuration of the actualization chain is the one Aristotle describes first. In \textit{De Motu Animalium} 7 he writes:

\begin{adjustwidth}{0.5in}{0in}
And so what we do without reflection, we do quickly. For when a man is actually using perception or imagination or thought in relation to that for the sake of which, what he desires he does at once. For the actualizing of desire is a substitute for inquiry or thinking. I want to drink, says appetite; this is drink, says sense or imagination or thought: straightaway I drink. In this way living creatures are impelled to move and to act, and desire is the last cause of movement, and desire arises through perception or through imagination and thought. And things that desire to act make and act sometimes from appetite or impulse and sometimes from wish (Aristotle, \textit{De Motu Animalium}, 701a29--b1).
\end{adjustwidth}

Three observations follow. First, ``the actualizing of desire is a substitute for inquiry or thinking'': appetitive action does not wait on deliberation, for the \textit{phantasma} that presents the object as drink supplies the functional equivalent of a conclusion, and the desire issues in motion before any evaluatively complex \textit{doxa} can engage. Second, the disjunction ``sense or imagination or thought'' names three channels by which the object may be recognized, but in the epithymetic case any one suffices---perception alone, without deliberation; the recognizing channel varies while the appetite-to-motion structure holds constant. Third, the framing sentence---``what we do without reflection, we do quickly''---presents this unreflective swiftness as the normal case for appetitive action, not as a deviation from a deliberative default.

Epithymetic action has a floor below even the pursuit of the pleasant: the avoidance of the bad in its most immediate, purely sensory form. The animal that startles at a sudden noise, the human being who flinches at an unexpected blow, recoils before any deliberative register is engaged. The genesis of such movement lies in the \textit{kritikon} of perception itself: in the very act of perceiving, in the `now', the discriminative dimension of \textit{aisthēsis} institutes the most basic good and bad, the to-be-pursued and the to-be-avoided (\textit{De Anima} III.7, 431a8--12), and from what it discriminates as bad the \textit{orexis} recoils. \textit{Pathos} here does not generate the movement; it is the very capability of being-moved---the being-affected (\textit{paschein}) that such discrimination and recoil presuppose. This is the one type that can fire pre-doxastically. It need not, however: epithymetic desire equally runs through a \textit{doxa}, as when one judges another to be a pleasing companion and pursues accordingly---an evaluative assent, answerable to how things are, and so capable of being mistaken. What individuates the type is its source in appetite, not the level of doxastic articulation at which it happens to operate.

\textbf{Type 2: Thymotic Action (\textit{thymos}).} Where appetite pursues the pleasant and shuns the painful, \textit{thymos} answers a slight, defends honour, and asserts the self---and it does so always in relation to another. Aristotle's anger arises definitionally from a perceived slight, an appraisal that presupposes the agent's standing, the offender's intention, and the social matrix within which the slight registers as undeserved (Aristotle, \textit{Rhetoric}, 1378a31--b5). \textit{Thymos} is in this sense irreducibly social: its object is disclosed within a \textit{toward-whom}, a Being-with (\textit{Mitsein}) in which others are met as those who can honour or wrong one. On Gross's reading the emotions ``tie humans in a unique way to their embodiment\ldots\ by determining the possibilities for moving about a shared world'' (Gross, ``Introduction,'' \textit{Heidegger and Rhetoric}, p.~26); and the Heidegger lecture Gross there translates puts the point in a formula---``the \textit{eidos} of the \textit{pathē} is a disposition toward other humans, a Being-in-the-world'' (Heidegger, quoted in Gross, p.~26). Anger is one such disposition. Because it answers an appraisal, thymotic action is doxastic throughout; and it admits of degree, Heidegger marking the distance between what is grasped ``in a state of arousal, in blind passion'' and what is grasped ``in clear, lucid resolution'' (\textit{BCAP} p.~99).

\textbf{Type 3: Bouletic Action (\textit{boulēsis}).} \textit{Boulēsis} is the rational wish for an end. Its object is the good held as end---``happiness and its constituents,'' which Aristotle names as the end at which every individual man and all men in common aim, and which determines what they choose and what they avoid (Aristotle, \textit{Rhetoric}, I.5, 1360b4--6). What marks \textit{boulēsis} off from \textit{prohairesis} is precisely that it is wish for the \textit{end} and not yet choice of the means: ``we wish to be healthy,'' Aristotle observes, ``but we choose the acts which will make us healthy'' (Aristotle, \textit{Nicomachean Ethics}, III.2, 1111b21--30). Because it reaches toward the end rather than instituting the realizable, \textit{boulēsis} may aim where \textit{prohairesis} cannot: there may be a wish even for impossibles, ``e.g.\ for immortality,'' or for what depends not on us but on others (Aristotle, \textit{Nicomachean Ethics}, III.2, 1111b21--30). Heidegger marks the same boundary---wishing, unlike resolute choice, ``can be directed at something that is impossible'' (\textit{BCAP} p.~99). The bouletic moment is thus the soul's orientation toward an end held as good, whether or not the path to it is yet in view: as when, moved to pity by another's suffering, one wishes the suffering relieved before any means has been weighed.

\textbf{Type 4: Prohairetic Action (\textit{prohairesis}).} \textit{Prohairesis} is the deliberate choice of the means by which a wished end is realized---deliberation completed in a resolution that institutes action. Aristotle's practical syllogism exhibits its form:

\begin{adjustwidth}{0.5in}{0in}
But how is it that thought is sometimes followed by action, sometimes not; sometimes by movement, sometimes not? What happens seems parallel to the case of thinking and inferring about the immovable objects. There the end is the truth seen (for, when one thinks the two propositions, one thinks and puts together the conclusion), but here the two propositions result in a conclusion which is an action---for example, whenever one thinks that every man ought to walk, and that one is a man oneself, straightaway one walks; or that, in this case, no man should walk, one is a man: straightaway one remains at rest. And one so acts in the two cases provided that there is nothing to compel or to prevent (Aristotle, \textit{De Motu Animalium}, 701a7--16).
\end{adjustwidth}

The conclusion here is not a proposition but the action itself---``straightaway one walks.'' What the syllogism stages is the passage from a held end to the realizable act that serves it. Deliberation, Aristotle insists, is not about ends but about what contributes to them: ``having set the end, they consider how and by what means it is to be attained'' (Aristotle, \textit{Nicomachean Ethics}, III.3, 1112b12--19), and it ranges only over ``things that are in our power and can be done'' (Aristotle, \textit{Nicomachean Ethics}, III.3, 1112a31). Hence \textit{prohairesis} is bounded in just the way \textit{boulēsis} is not. It goes after the \textit{praktón}, ``that which is decisive for a concern in the moment'' (\textit{BCAP} p.~98); and where wish may reach toward the impossible, choice ``never goes after something that is impossible'' but ``always goes after something that is under our control,'' carrying through to the \textit{éschaton}---the point ``that I grasp, that I genuinely institute through action'' (\textit{BCAP} p.~99). Deliberate choice can also sediment into readiness: the expert pianist no longer weighs each fingering afresh but realizes a passage through means long since chosen and now habitual---prohairetic resolution compressed by practice into fluency. How such sedimentation cultivates a disposition, and how that disposition differs when the end lies in a product rather than in the doing itself, belongs to the bivalence of \textit{hexis} taken up below.

Across all four types the division is by source-desire alone; it is orthogonal to a second distinction---whether the action's end lies beyond it, in a product (\textit{poiēsis}), or is immanent in the doing (\textit{praxis})---which cross-cuts the four and is the matter of the \textit{hexis}-bivalence below. And within each of the doxastic types, what decides whether the desire concretizes into a determinate civic passion---anger, fear, pity, shame---is whether the \textit{doxa} ratifies a content that is evaluatively complex. That condition is what the next subsection makes precise.

%---END 04-FourTypes---

%---BEGIN 05-DoxRatification (target 700w)---
\subsubsection*{Doxastic Ratification and Its Content-Condition}

Higher-order \textit{pathē} arise only when two jointly necessary conditions obtain: first, \textit{doxa} must form as an involuntary \textit{taking-as-true}, a propositional commitment that requires \textit{logos}, \textit{pistis}, and \textit{pepoithēsis} and is accordingly exclusive to rational beings (Aristotle, \textit{De Anima}, 428a19--24); second, the content of that \textit{doxa} must engage evaluatively complex categories---slight, threat, undeserved suffering, disgrace, or their kin. We may call \textit{doxa}'s affirmative role here \textit{doxastic ratification}, and the requirement that what it ratifies be evaluatively complex its \textit{content-conditionality}. Neither condition alone suffices. A \textit{doxa} whose content is evaluatively inert---``every man ought to walk; I am a man''---issues in action without higher-order affect (Aristotle, \textit{De Motu Animalium}, 701a7--16). Conversely, a bare \textit{phantasma} of something fearful, however vivid, leaves the soul in the condition of one who merely contemplates a painted terror, unmoved in the register of genuine \textit{pathos}. This conjunction is thus not a single switch but a two-dimensional threshold: \textit{doxa} is the necessary rational axis, evaluative complexity is the necessary material axis, and only their intersection produces the determinate higher-order \textit{pathē} that Aristotle catalogues across \textit{Rhetoric} II.

What such \textit{doxa}-formation minimally involves can be stated precisely. A \textit{doxa} takes an appearance as bearing on reality; as an assent answerable to how things are, it turns false not through any act of the believer but when the facts change unremarked. Nor is the assent at the believer's discretion. One may voluntarily picture two courses of action and their outcomes---this is deliberative \textit{phantasia} at work---but a \textit{doxa} forms only when one settles, involuntarily, on the course that is better: a judgment reached through comparison, not an image summoned at will. This is why we cannot believe whatever we wish, but only what we take to be true. \textit{Doxa}, in short, is the soul's affirmative ratification of an appearance---a propositional \textit{taking-as-true}, not a disposition---held as true until reality or further reflection revises it.

The textual warrant for this structure appears most concisely at \textit{De Anima} III.3, 427b21--24, where Aristotle contrasts mere imaginative presentation with doxastic ratification. When the soul merely imagines fearful things, it remains ``unaffected''---the \textit{phantasma} operates, but no higher-order \textit{pathe} crystallizes; yet when \textit{doxa} propositionally ratifies the aspectual presentation, we are affected \textit{euthys}, immediately (Aristotle, \textit{De Anima}, 427b21--24). The adverb is decisive: it marks the \textit{pathe} not as a gradual accretion but as an abrupt qualitative shift, a threshold-crossing that occurs the instant \textit{doxa} engages the evaluatively charged content of the \textit{phantasma}. Indeed, the immediacy (\textit{euthys}) signals that ratification is not a deliberative procedure but a structural feature of the actualization chain at $A_3$: when the two conditions converge, the chain passes through evaluative \textit{articulational concretion}; when they do not, it bypasses that register entirely.

This content-conditional dimension receives its most explicit articulation in the \textit{Rhetoric}. Aristotle defines \textit{pathē} as those feelings that ``so change men as to affect their judgements,'' coupling affective transformation with evaluative-judgemental content (Aristotle, \textit{Rhetoric}, 1378a21--22). The definition is not incidental; it specifies that the judgement in question must bear upon matters of evaluative weight---matters, that is, whose propositional content engages categories such as justice, honour, or shame. Anger furnishes the paradigmatic instance: it arises definitionally from a perceived slight, an evaluatively complex category that presupposes the perceiver's standing, the slighter's intention, and the social matrix within which the slight registers as undeserved (Aristotle, \textit{Rhetoric}, 1378a31--b5). Similarly, fear requires not merely the \textit{phantasma} of something destructive but the \textit{doxa} that a destructive evil is genuinely impending and pertinent to the one who fears (Aristotle, \textit{Rhetoric}, 1382a21--22). Thus this content-conditionality is what determines, across the four types of action, whether higher-order \textit{pathē} concretize. Epithymetic action alone can run below this threshold: in its pre-doxastic case the \textit{phantasma} drives the chain from appetite through organic preparation to movement without doxastic ratification on any axis (Aristotle, \textit{De Motu Animalium}, 702a17--19). Wherever \textit{doxa} does form---in epithymetic action's doxastic case, and throughout thymotic, bouletic, and prohairetic action---a second condition still governs the outcome: when the ratified content is evaluatively inert, as in the procedural \textit{doxa} that concludes a deliberation about means, no higher-order \textit{pathos} arises; when \textit{doxa} ratifies an evaluatively complex \textit{phantasma}, the chain runs through the full \textit{articulational concretion} that yields anger, fear, pity, shame, or their cognates.

A concrete instance shows what turns on the content \textit{doxa} ratifies. Walking at night down an unfamiliar alley, one is pervaded by a vague anxiety---a diffuse attunement in which the surroundings disclose themselves as faintly menacing, though nothing in particular is yet feared. The latent orectic charge in play---the \textit{resonant epithymia} deposited in perception---is not itself felt in the moment; what is felt is only this indeterminate unease. Then a hooded figure springs from behind a trash can. The instant \textit{doxa} ratifies the \textit{phantasma} as threatening, the diffuse anxiety crystallizes into determinate fear---the judgement of imminent harm, the spike of dread, the impulse to flee, the pulse that readies the body. Had the figure been taken for a friend, or for a windblown bag, the unease would have subsided rather than concretized; what converts a vague anxiety into a determinate \textit{pathos} is solely what \textit{doxa} takes the appearance to be.

The relation runs in the other direction as well. When \textit{doxa} ratifies an evaluatively complex \textit{phantasma}, the opinion that crystallizes does not stand free of the affect within which the whole process is situated; it inherits its coloring. The judgment that the beacon signals danger is no affectively neutral proposition; it is saturated with the fear the perception aroused---a charge co-present in the content the \textit{doxa} ratifies---and this saturation conditions both the urgency and the character of the ensuing action. The point is unsurprising once we recall that the very feelings that alter judgment are themselves ``attended by pain or pleasure'' (Aristotle, \textit{Rhetoric}, 1378a20--23): affect does not arrive after the opinion to disturb it but belongs to the content the opinion takes up. Doxastic ratification, accordingly, does not purify the rational moment of its \textit{pathetic} origins; it translates affective disclosure into evaluable content without extinguishing the affect---which is why emotion and mood can shape what we take to be true.

This structural account illuminates the therapeutic and rhetorical possibilities latent in the symmetry of this conditionality. A rhetor who deploys \textit{enargeia}---vivid, concrete, affectively charged imagery---can perhaps induce an evaluatively complex \textit{phantasma} that transforms a routine, affectively inert encounter into a charged one in which higher-order \textit{pathē} concretize. The same mechanism runs in both directions. If \textit{enargeia} can supply a \textit{phantasma} evaluatively charged enough for \textit{doxa} to ratify as a slight or a threat, reframing can as readily defuse it: recast the apparent slight as an innocent oversight, and the \textit{phantasma} \textit{doxa} now assents to is evaluatively inert, returning the audience to composure. The process is in this sense symmetrical: one and the same threshold is crossed whether affect is kindled or dissolved, and what decides the outcome is not the direction of crossing but whether the content to which \textit{doxa} assents is evaluatively complex. Von Uexküll's account of the child at play furnishes a vivid instance. A child playing ``the story of Hansel and Gretel, the witch, and the gingerbread house'' becomes so gripped by the imagined witch that she cries, ``Get the witch out of here; I can't stand to see her repulsive face any more!'' (von Uexküll, \textit{A Foray into the Worlds of Animals and Humans}, pp.~122--126). The \textit{phantasma} has become vivid enough to grip the child affectively despite her standing \textit{doxa} that the scene is merely play; the affective response is keyed to the appearance, not to the judgment she could still avow, and for the moment the \textit{doxa} that had governed the activity no longer holds her. The episode shows at once that a \textit{doxa} is held only so long as it continues to convince---held-as-true until revised---and that a sufficiently charged \textit{phantasma} can override a standing one. Struever, reading Heidegger's account of the \textit{pathē}, places the moment of their force at the point of articulated judgement: passions ``intrude on logos---here articulated judgements'' (Struever, \textit{Heidegger and Rhetoric}, p.~109). The passions do not float free of rational structure; they crystallize at the juncture where propositional commitment meets evaluative substance---that is, where \textit{doxa} ratifies a content that is evaluatively complex.

Doxastic ratification thus occupies the pivotal position within the actualization chain at $A_3$. It is what determines whether the chain's passage through completed cognitive actuality includes evaluative articulation or bypasses it; whether the \textit{phantasma} stabilized at $A_2$ becomes the nucleus of a determinate higher-order \textit{pathe} or remains, so to speak, affectively quiescent. Yet ratification does not operate in a vacuum. Whether the content it affirms is evaluatively complex is itself conditioned by the dispositional ground upon which \textit{doxa} forms---by the held \textit{doxai} and cultivated \textit{hexeis} that function as additional unmoved originators feeding the motion toward $A_3$ (Aristotle, \textit{Categories}, 8b25--9a13). How that dispositional ground is itself structured---whether as the transparent routine of \textit{technē} or as the open, \textit{kairos}-responsive\footnote{\textit{Kairos} (\gk{καιρός}) is the opportune or decisive moment---the qualitatively right time for action, as against \textit{chronos}, time as measured duration. A \textit{kairos}-responsive disposition is cultivated not to apply a fixed rule but to meet the singular demand of the situation as it arises---the readiness Heidegger glosses as ``holding-oneself-open'' rather than routine (\textit{BCAP}, p.~128).} composure of the ethical \textit{hexis}---is accordingly the question that this analysis now makes urgent.

%---END 05-DoxRatification---

%---BEGIN 06-HexisBivalence (target 900w)---
\subsubsection*{\textit{Hexis} Bivalence: \textit{Technē} and Ethical Hexis}

\textit{Hexis} names a stable, settled disposition encompassing virtues, crafts, knowledge, and bodily states alike, marked by its durability and resistance to easy displacement, such that knowledge and moral excellence count as paradigmatic \textit{hexeis} while passing moods and fleeting opinions do not (\textit{Categories} 8, 8b25--9a13). The dispositional ground that conditions which type of action an agent enters is, however, itself bivalent---and it is this bivalence, articulated through Heidegger's reading of Aristotle in the 1924 lectures, that proves decisive for understanding how the totality of involvements constituting practical life is structured by two fundamentally different modes of cultivated readiness. Aristotle defines \textit{hexis} more broadly in the \textit{Metaphysics} as the activity of haver and had, a disposition toward being well or ill disposed (\textit{Metaphysics} $\Delta$.20, 1022b4). Furthermore, in the \textit{Nicomachean Ethics} he insists that moral excellence comes about not by nature but as a result of habit: we become just by doing just acts, temperate by doing temperate acts (\textit{Nicomachean Ethics} II.1, 1103a14--b25). The universal premises of the practical syllogism---as when one holds that ``every man ought to walk'' and recognizes oneself as a man (\textit{De Motu Animalium} 7, 701a7--16)---are \textit{held doxai} that an agent's \textit{hexeis} may incline him to hold; they are not themselves \textit{hexeis}, for a \textit{doxa} is a taking-as-true, not a disposition. \textit{Hexis} proper is the broader dispositional ground within which bodily, affective, and practical-rational states belong. Indeed, Heidegger identifies this breadth as essential to understanding \textit{pathē}: ``\textit{Hexis} is nothing other than a how of \textit{pathos}, being-out-of-composure, in relation to being-composed-as-to \ldots'' (BCAP p.~125). Thus \textit{hexis} is not merely a cognitive deposit but the dispositional form through which \textit{pathē} themselves become determinate: one and the same \textit{pathos}---fear, say---takes one shape in the coward and another in the courageous person, since the \textit{hexis} is the cultivated \textit{how} of standing in that \textit{pathos}, as courage is the cultivated comportment toward fear (\textit{Nicomachean Ethics} III.6--7). It is along this axis that the bivalent reading emerges.

Heidegger's lectures on Aristotle's \textit{Rhetoric} develop a first pole of this bivalence under the heading of \textit{technē}---the settled disposition cultivated through training in a productive art. The structural logic is ergonomic: training progressively reduces the deliberative burden that accompanies unfamiliar operations, until the practitioner's orientation to the task becomes transparent, procedural, and effectively automatic---transparent in just the way Heidegger's ready-to-hand equipment is, the hammer withdrawing from explicit notice into the hammering itself and becoming conspicuous only when it fails (\textit{Being and Time}, \S15, H.69; Eng.~98). Heidegger states:

\begin{adjustwidth}{0.5in}{0in}
``Through practice, by frequently-undergoing, it comes about that being-oriented puts the prescription further and further out of play. Training has the precise sense of reducing deliberation insofar as it is through training that the completedness of attaining a result comes about'' (BCAP p.~127).
\end{adjustwidth}

The expert pianist's hands, the carpenter's chisel-grip, the orator's settled cadence---each exemplifies a \textit{technē} that fires the chain of action through transparent procedural \textit{doxa} without engaging the evaluative axis of higher-order \textit{pathē}. The ``completedness of attaining a result'' is the defining structure here: the result is external to the activity, the operation is means-to-end, and training cultivates routine. Accordingly, when the operative \textit{doxa} is a held procedural one---a craft-directive drawn upon rather than freshly weighed---the chain runs forward from $A_2$ through $A_3$ to $A_4$ without fresh evaluative ratification. The \textit{doxa} is present, but it operates as a standing directive rather than as a living encounter with the particular situation's demand. Aristotle's observation that ``movement is a kind of activity---an imperfect kind'' applies here with particular force: \textit{technē} reduces the imperfection not by completing the motion into full \textit{energeia} but by abbreviating it, cutting short the deliberative arc that would otherwise hold the agent open to the situation's complexity (Aristotle, \textit{De Anima} II.5, 417a14--20).

The second pole---the ethical \textit{hexis}---names a fundamentally different dispositional structure, one that Heidegger contrasts sharply with the routine logic of \textit{technē}:

\begin{adjustwidth}{0.5in}{0in}
``Cultivating \textit{hexis} never depends on an operation, a routine. In an operation, the moment is destroyed. Every completedness, as settled routine, breaks down in the face of the moment. Appropriation and cultivation of \textit{hexis} through habituation means nothing other than correct repetition\ldots. The distinction lies in the fact that \textit{praxis} depends on the \textit{how}. The how is only appropriated in such a way that the human being enables himself \textit{to be composed at each moment}; not routine but holding-oneself-open, \textit{dynamis} in the \textit{mesōtes}'' (BCAP p.~128).
\end{adjustwidth}

The ethical \textit{hexis}, then, is the cultivation of a readiness that does not destroy the moment (\textit{Augenblick}) but sustains it; it is the \textit{hexis} through which an agent meets a situation in fresh evaluative resolution---neither bypassing \textit{doxa}, as appetite may, nor running it on routine, but holding oneself open for fresh \textit{kairos}-resolution. Heidegger is explicit that there can be no \textit{technē} adequate to this demand: ``this determination should not be conceived as though there were a \textit{technē} for this taking-opportunities and venturing-out into the \textit{deina} of life'' (BCAP pp.~122--123). The \textit{kairos} is precisely what \textit{technē} cannot resolve, for its routine structure forecloses the very openness that the singular moment demands. Similarly, the canonical practical-rational \textit{hexis} that exemplifies this holding-oneself-open is \textit{phronēsis}---``a true and reasoned state of capacity to act with regard to the things that are good or bad for man''---which Aristotle distinguishes from \textit{technē} on the ground that \textit{praxis} has no external \textit{ergon} but is its own end (\textit{Nicomachean Ethics} VI.5, 1140a24--b30).

The receptive character shared by Aristotelian \textit{aisthēsis} and Heideggerian \textit{Befindlichkeit} proves illuminating here, for both structures disclose the agent as already affected, already disposed, before any deliberative act begins (Aristotle, \textit{De Anima} II.5, 416b33--417a1; Heidegger, \textit{Being and Time}, \S 29, H.137; Eng.~p.~176). The ethical \textit{hexis} uniquely exercises a two-level reach across this receptive ground: it modulates both the basic-valence input at $A_1$---the \textit{Stimmung}-saturation of perceptual disclosure, in which the world is always already encountered as mattering in some determinate way---and the orthogonal dimension of \textit{doxa} at $A_3$, where the soul's readiness to ratify the \textit{phantasma} stabilized at $A_2$ as evaluatively charged determines whether higher-order \textit{pathē} arise. \textit{Technē}, by contrast, runs the chain forward without engaging the evaluative axis at all; hence it neither deepens nor transforms the affective ground from which perception arises. Of the \textit{pathē} more broadly, Struever observes---reading Heidegger's 1924 lectures---that passion ``is both motion and a cause of motions; it is a capacity for altering,'' so that ``this is where passions intrude on logos---here articulated judgements'' (Struever, ``Alltäglichkeit, Timefulness, in the Heideggerian Program,'' in \textit{Heidegger and Rhetoric}, p.~109). The ethical \textit{hexis} is thus the cultivation through which the body's dispositional condition becomes ready for fresh evaluative resolution, holding itself open to the concrete situation rather than abbreviating it through routine.

The bivalence is nowhere clearer than in a single practitioner exercising both dispositions. In suturing a wound the surgeon exercises \textit{technē}: the end is a detachable product---the closed wound---and the work is judged instrumentally, by the efficiency and quality with which that external result is attained. In sitting with a dying patient the same surgeon exercises an ethical \textit{hexis}: here there is no detachable artifact, only the conduct itself---speaking truthfully, courageously, compassionately---in which the excellence wholly consists. The end of the first lies beyond the doing, in the work; the end of the second is immanent in the doing. And because an ethical \textit{hexis} is the cultivated \textit{how} of standing in a \textit{pathos}, it is named for the comportment it achieves, not the affect it presupposes: we call the surgeon courageous, not fearful, though fear is the very condition courage masters (\textit{Nicomachean Ethics} III.6--7).

This teleological difference is, at bottom, the difference between two kinds of actualization. \textit{Poiēsis}, the making that \textit{technē} governs, is a \textit{kinēsis}: its end lies in a product beyond the activity, so the activity is incomplete in itself. \textit{Praxis}, the doing that ethical \textit{hexis} governs, is an \textit{energeia}: its end is internal, so that one acts well and has acted well at once. Aristotle marks the contrast by a tense-test: ``it is not true that at the same time we are \ldots\ building and have built \ldots; but it is the same thing that at the same time has seen and is seeing \ldots. The latter sort of process, then, I call an actuality, and the former a movement'' (\textit{Metaphysics} IX.6, 1048b30--36). The bivalence of \textit{hexis} thus rests on the same ontological seam as the distinction between motion and actuality itself: a \textit{kinēsis} is an \textit{energeia atelēs}, an incomplete actuality, and \textit{technē} is the disposition perfected in such incomplete actualization, while ethical \textit{hexis} is the disposition perfected in the complete actualization that is \textit{eupraxia}.\footnote{This maps the standard \textit{Metaphysics} IX.6 contrast onto the \textit{praxis}/\textit{poiēsis} pair as a structural homology rather than a strict identity; Ackrill's puzzles about whether every \textit{praxis} is an \textit{energeia}---some actions appear to leave results that outlast them---caution against a clean identification.}

This bivalence is not a further kind of action set beside the four but a distinction that cross-cuts them: any of the four may terminate either in a making, whose end is a product beyond the doing, or in a doing whose end is internal to it. What varies is then not the source-desire but the dispositional mode through which the chain runs. Appetite that fires pre-doxastically engages no settled disposition at all, running from $A_1$ through $A_2$ to $A_4$ without dispositional mediation. Where the chain runs through \textit{technē} (the \textit{hexis} of production, \textit{poiēsis}), the cultivated craft-recognition fires the chain on a trained judgment, its passage through $A_3$ effectively transparent. And where the agent meets the situation in fresh evaluative resolution, it routes through the ethical \textit{hexis}, \textit{doxa} ratifying the \textit{phantasma} in a genuine encounter with the particular \textit{kairos}. Subsequently, every completed action at $A_4$ sediments back into the agent's \textit{hexeis}, the kind sedimented keyed to the terminus: production deepens \textit{technē}, reducing deliberation and consolidating routine, while a doing whose end is its own performance deepens the ethical \textit{hexis}, strengthening the agent's capacity for fresh evaluative openness. It is this diachronic sedimentation---whereby the dispositional ground is itself reshaped by the actions it enables---that constitutes the recursive structure of the \textit{pathetic} loop.

%---END 06-HexisBivalence---

%---BEGIN 07-PatheticLoop (target 700w)---
\subsubsection*{The \textit{Pathetic} Loop: Synchronic Pass and Diachronic Sedimentation}

The \textit{pathetic} loop operates in two temporal registers---synchronic and diachronic---and the relationship between them is what gives the actualization chain its recursive, spiraling character. In its synchronic register, the loop describes a single, directional chain-pass from $A_0$ through $A_4$ within one discrete event: the dyadic encounter of sensible object and \textit{aisthētikon} yields completed perception at $A_1$, whose dual-resonance residues are stabilized as \textit{phantasma} at $A_2$, oriented cognitively at $A_3$, and consummated in \textit{praxis} at $A_4$. Within this pass, whether \textit{doxa} ratifies an evaluatively complex content at $A_3$ determines whether the \textit{phantasma} produces higher-order \textit{pathē}; and when it does, the resultant affective-orectic mobilization converts generic \textit{orexis} into species-specific desire---anger, fear, pity, shame---that drives $M_3 \rightarrow A_4$. This synchronic register is what the \textit{Rhetoric}'s canonical definition of \textit{pathē} articulates: feelings ``that so change men as to affect their judgements'' (Aristotle, \textit{Rhetoric}, 1378a20--22). The judgement-modifying function is not an accidental feature of emotion but the very structure of the synchronic pass at the $A_3$ node, where \textit{doxa} and \textit{pathe} meet in a single operation.

The diachronic register, by contrast, describes what happens \textit{across} chain-passes. Residues of prior actualizations do not simply vanish upon the completion of $A_4$; they sediment as dispositional ground---as \textit{hexis}---modulating subsequent passes through two distinct mechanisms. The first is what we may call \textit{Stimmung}-saturation: the agent's accumulated affective history biases basic valence at $A_1$, coloring the perceptual disclosure of present objects before any deliberative operation can intervene. Aristotle provides the canonical description of this mechanism in \textit{De Insomniis}: ``even when the external object of perception has departed, the impressions it has made persist, and are themselves objects of perception,'' and under the influence of emotion ``the coward when excited by fear'' thinks he sees his foes approaching, while ``the amorous person'' sees the object of desire, ``and the more deeply one is under the influence of the emotion, the less similarity is required to give rise to these impressions'' (Aristotle, \textit{De Insomniis}, 460b1--10). What Aristotle specifies here is not a marginal pathological case but rather the general structure of affective priming: the dispositional ground laid down by prior chain-passes saturates the receptive capacity of \textit{aisthēsis} itself, so that the perceptual field is always already tonally inflected. Heidegger names this inflection \textit{Befindlichkeit}, the existentiale through which Dasein finds itself always already disposed in a world that is disclosed not neutrally but atmospherically, as \textit{Stimmung} (Heidegger, \textit{Being and Time}, \S 29, H.137; Eng.~p.~176). The receptive character shared by \textit{aisthēsis} and \textit{Befindlichkeit} is thus not merely analogical; both name the structure whereby an entity is affected by beings prior to any thematic judgement, a structure that is itself constitutively shaped by what has come before.

The second diachronic mechanism is corrective-faculty impairment. The same \textit{De Insomniis} passage specifies that existing emotional states modulate doxastic ratification at $A_2 \rightarrow A_3$: the corrective capacity that would normally resist the \textit{taking-as-true} of evaluatively distorting \textit{phantasmata} is weakened under emotional excitement, so that \textit{doxa} ratifies more readily the images consonant with the agent's current affective register (Aristotle, \textit{De Insomniis}, 459a24--28; 460b3--16). Hence neither mechanism introduces new transitions into the chain; both bias existing ones. Heidegger's reading of \textit{hexis} in \textit{Basic Concepts of Aristotelian Philosophy} confirms this structural point: \textit{hexis} is ``nothing other than a how of \textit{pathos}, being-out-of-composure, in relation to being-composed-as-to\ldots'' (BCAP p.~125). The dispositional ground is not a container for emotions but the very manner in which the agent is susceptible to being affected---the \textit{how} of \textit{paschein} rather than its \textit{what}.

Indeed, the \textit{Rhetoric}'s judgement-modifying definition and the \textit{De Insomniis} mechanism are not two separate doctrines but two articulations of one loop-structure operating across two temporal registers. Synchronically, emotion modifies the very judgement that would otherwise check it; diachronically, the sedimented residue of that modification reshapes the perceptual and doxastic conditions for every subsequent pass. Accordingly, the chain is not merely iterative but spiraling: $A_4$ completion deposits fresh \textit{hexis}-formation, and the next chain-pass enters at a new $A_0$ with the dispositional ground subtly---or perhaps drastically---shifted. The agent who has just acted in anger is not the same agent who will next encounter a perceived slight, because the act itself has reshaped the affective architecture through which the world is disclosed. Thus passions ``alter judgements'' precisely where they ``intrude on \textit{logos}'' (Struever, \textit{Heidegger and Rhetoric}, p.~109), and this intrusion is not an episodic disruption but the ongoing self-modification of the dispositional ground. Gross captures the anti-internalist force of this thesis: ``Heidegger characterizes \textit{pathos} (variously `passion,' `affect,' `mood,' or `emotion') as the very condition for the possibility of rational discourse, or \textit{logos}. No cynical and crowd-pleasing addition to \textit{logos}, \textit{pathos} is the very substance in which propositional thought finds its objects and its motivation'' (Gross, ``Introduction,'' in \textit{Heidegger and Rhetoric}, 2005, p.~4). \textit{Pathos} is not added to a neutral deliberative process; it constitutes the medium in which deliberation occurs at all.

This recursion has a precise mechanism. As Nussbaum observes, translating the \textit{Nicomachean Ethics}, ``All men strive for the apparent good (\textit{phainomenon agathon}): but no one is in control of the appearing (\textit{phantasia}): the way the end appears (\textit{phainetai}) to someone depends on what sort of a man he is'' (Nussbaum, ``The Role of \textit{Phantasia} in Aristotle's Explanation of Action,'' pp.~31--33). The sort of man one is---one's sedimented \textit{hexis}---conditions the very \textit{phantasia} through which the good presents itself, and so conditions the content available for doxastic ratification. The loop is therefore no simple feedback mechanism but a genuinely recursive structure: disposition shapes appearance, appearance shapes action, and action, through its \textit{pathetic} residue, reshapes disposition. This is what makes ethical transformation so difficult on Aristotle's account: a vicious \textit{hexis} reproduces itself by conditioning the very appearances available to the agent, narrowing the content that might otherwise present a more adequate good, so that the loop, once established, runs with a momentum not easily interrupted from within.

A personal example makes the appetitive route concrete. Having given up smoking after fifteen years, I find that long after the physiological withdrawal has passed the craving persists: the cigarette still presents itself as desirable even as I judge, without wavering, that it is harmful. Years of indulgence sedimented an appetitive \textit{hexis}---an intemperance, in Aristotle's terms (\textit{Nicomachean Ethics} III.10--12)---that continues to condition the \textit{phantasia} through which the object appears, exactly as Nussbaum's gloss predicts: the appearing is not in my control. Abstention slowly cultivates the contrary disposition, temperance, but the entrenched \textit{hexis} does not dissolve at once, and its persistence is just the self-reinforcing momentum of the loop. This is the appetitive analogue of the graded feedback that distinguishes the other routes: as \textit{poiēsis} sediments \textit{technē} and \textit{praxis} sediments virtue, so repeated appetitive action sediments \textit{empeiría} (\textit{Metaphysics} A.1, 980b28--981a4) together with the temperance or intemperance that disposes one toward the body's pleasures. No completed action is disposition-neutral; each leaves its residue in the kind of agent one is becoming.

What Heidegger names \textit{Geschichtlichkeit}---Dasein's historicity---is, at its deepest structural level, this diachronic register of the \textit{pathetic} loop's recursive self-modification. ``As historical, Dasein is possible only by reason of its temporality, and temporality temporalizes itself in the ecstatico-horizonal unity of its raptures'' (Heidegger, \textit{Being and Time}, \S 74, H.385--386; Eng.~p.~437). The Heideggerian articulation does not supplement the Aristotelian structure; it names what the Aristotelian text already specifies---that the totality of involvements through which \textit{orexis} finds its objects is never fixed but perpetually reconstituted by the sedimented history of the agent's own affective life. The agent's dispositional ground is, accordingly, not a static substrate but a living, self-revising archive---the diachronic dimension of the \textit{pathetic} loop spiraling forward into each new encounter with the world.

%---END 07-PatheticLoop---

%---BEGIN 08-Synthesis (target 700w)---
\subsubsection*{Synthesis: Emotion as the Affective Architecture of Being-in-the-World}

Emotion is not an add-on to rationality, not a disruption of cognition; it is the way the soul's self-movement discloses significance at the precise point where significance must become motion. What \textit{phantasia} presents, \textit{doxa} commits to; what \textit{doxa} commits to, emotion mobilizes; what emotion mobilizes, action completes. Each moment in this sequence is \textit{kinēsis}, and each is \textit{kinēsis} of an ensouled being whose very being is to-be-moved. Aristotle insists in \textit{De Anima} I.4 that ``being pained or pleased, or thinking'' are not activities of the soul in isolation but movements of the composite, ensouled being as such (Aristotle, \textit{De Anima}, 408b5--7). Hence the actualization chain traced across this analysis---from the joint actuality of sensible object and \textit{aisthētikon} at $A_0$, through the dual-resonance deposit at $A_1$, to the stabilized \textit{phantasma} at $A_2$, the doxastically ratified cognitive actuality at $A_3$, and the completed \textit{praxis} at $A_4$---is not a theoretical model imposed on Aristotle's psychology from without but the structural articulation of how rhetorical-affective being unfolds. At every node, it is the soul's \textit{pathetic} dimension---its capacity for being affected---\textit{paschein} in its preservative, \textit{sōtēria} sense, the being-affected that fulfils a power rather than destroying it---that sustains the chain's momentum. Indeed, Aristotle characterizes perception and appetite as ``one in substrate, different in being'' (\textit{De Anima} III.7, 431a8--14), and the \textit{Rhetoric}'s canonical definition of the \textit{pathē} as judgement-modifying dispositions (\textit{Rhetoric} II.1, 1378a21--22) confirms that emotion operates not alongside but \textit{within} the very structure of evaluative disclosure. The soul does not first perceive, then feel, then judge; perception is already affectively saturated, and judgement is already orectically charged.

It is precisely this insight that Heidegger radicalizes in his reading of Aristotle's \textit{pathē}. In the \textit{Basic Concepts of Aristotelian Philosophy}, Heidegger describes the fundamental character of \textit{paschein} as ``affected by beings as world; \textit{paschein} as \textit{dechesthai}, `to perceive' and that primarily as \textit{pathos}, `becoming-affected,' `disposition,' being related to the world as in it, being in dealings with it'' (Heidegger, \textit{Basic Concepts of Aristotelian Philosophy}, pp.~277--280). The \textit{pathē}, on this account, are ``not psychic experiences and not in consciousness'' but ``a being-taken of human beings \textit{in their full being-in-the-world}'' (Heidegger, \textit{Basic Concepts of Aristotelian Philosophy}, pp.~133--134). Thus Heidegger refuses every internalist reduction: emotion is not a private episode occurring behind the eyes but a mode of world-disclosure, a way in which the \textit{totality of involvements} that constitutes a life becomes manifest to the being who lives it. The receptive character that Aristotle ascribes to \textit{aisthēsis}---the capacity to take on form without matter---finds its ontological correlate in what Heidegger calls \textit{Befindlichkeit}, the primordial disclosedness through which Dasein is always already delivered over to its situation (Heidegger, \textit{Being and Time}, BT \S29, H.134; Eng.~p.~174). Similarly, \textit{Stimmung} is not a subjective coloring layered over neutral facts; Heidegger insists that ``a mood assails us. It comes neither from `outside' nor from `inside', but arises out of Being-in-the-world, as a way of such Being'' (Heidegger, \textit{Being and Time}, BT \S29, H.137; Eng.~p.~176). The structural homology is thus precise: for Aristotle, \textit{pathos} is the kinetic-affective-evaluative dimension of how the world appears to us, grounded in \textit{phantasia}, shaped through involuntary \textit{doxa}, felt as pleasure or pain, and expressed as \textit{orexis} toward or away from what the soul takes to be significant; for Heidegger, \textit{Stimmung} is the ontological-atmospheric correlate, the attunement that constitutes the clearing in which beings can appear at all.

Accordingly, the convergence of Aristotelian and Heideggerian accounts is not merely historical but structural. Gross captures this convergence with particular force. Writing against the cognitivist reduction of emotion to neural mechanism, Gross argues that ``the tradition after Aristotle's \textit{Rhetoric} offers alternatives for understanding emotions as social phenomena, which is something that leading humanists like Martha Nussbaum and Richard Sorabji should have remembered as they rushed toward the latest brain science'' (Gross, \textit{Uncomfortable Situations}, p.~3). Subsequently, Gross articulates the positive implication of this recovery: ``Rhetoric, as a humanistic form of affordance theory, posits a situation that is not neutral but is rather `persuasive,' threatening and promising with respect to our human being-in-the-world'' (Gross, \textit{Uncomfortable Situations}, p.~20). It is this formulation---rhetoric as humanistic affordance theory---that maps directly onto the dissertation's synthesis. If the situation is never neutral but always already ``persuasive, threatening and promising,'' then emotion is not a contingent addition to an otherwise dispassionate encounter with the world; it is the medium through which the world's significance becomes actionable. Furthermore, as Struever shows in reading Heidegger on \textit{Rhetoric} II, passion ``is both motion and a cause of motions; it is a capacity for altering'' (Struever, \textit{Heidegger and Rhetoric}, p.~109). A constituted \textit{pathos} is in this sense doubly kinetic: it is itself a being-moved and, once established, a source of further motion (\textit{archē kinēseōs}, a principle of movement)---not the originator of the synchronic chain, which remains the realizable good working through \textit{orexis}, but the impetus that carries the diachronic loop from $A_4$ back through \textit{hexis}-sedimentation into the dispositional ground of the next perceptual encounter.

We may now state the closing formulation explicitly: emotion, in its full Aristotelian sense and in its Heideggerian radicalization, is nothing less than the \textbf{affective architecture of being-in-the-world}---the kinetic, evaluative, dispositional structure through which the soul discloses significance, mobilizes action, and sediments the traces of its own history into the ground from which future disclosure arises.

%---END 08-Synthesis---

%---BEGIN 09-DevNote (target 400w)---
\subsection*{Development Note}

Five questions raised by the foregoing analyses remain open for development in subsequent chapters; each marks a point at which the actualization chain's internal logic presses beyond the resources assembled thus far and accordingly demands further architectural elaboration. \begin{enumerate}

\item \textbf{The \textit{technē} / ethical-\textit{hexis} question for a single agent.} The bivalent reading of \textit{hexis} developed from Heidegger's treatment in the \textit{Basic Concepts of Aristotelian Philosophy} operates per-action-type: training reduces deliberation toward transparent routine in one register, while holding-oneself-open for fresh \textit{kairos}-resolution characterizes the other (Heidegger, \textit{Basic Concepts of Aristotelian Philosophy}, p.~128). Yet an individual agent may inhabit both registers simultaneously across different domains---the carpenter who is also a parent, the orator who is also a friend. Whether \textit{technē} and the ethical \textit{hexis} are mutually exclusive within a single life, or whether the diachronic chain layers heterogeneous dispositional strata across an agent's totality of involvements, is a question deferred to the capstone synthesis. The per-agent topology remains, for now, an open problem. \item \textbf{The \textit{Stimmung}-saturation deferral.} The two diachronic mechanisms identified---\textit{Stimmung}-saturation at $A_1$ and corrective-faculty impairment via accumulated \textit{phantasmata}---may or may not prove reducible to one another. Whether \textit{Stimmung}-saturation is genuinely a direct effect of the perceptual \textit{hexis}-state or whether it is perhaps a derivative biasing mediated through $A_2$ \textit{phantasma}-content cannot be resolved without the ontological resources of \textit{Befindlichkeit} as Heidegger develops it---the equiprimordial existentiale disclosing Dasein's thrownness prior to cognition (Heidegger, \textit{Being and Time}, \S29, H.~134; Eng.~p.~174). The present analyses treat the two mechanisms as distinct; the reduction-question is deferred. \item \textbf{Rhetoric and external situation modulation.} The \textit{pathetic} loop's diachronic recursion, as articulated here, operates largely from within the agent's own dispositional ground. How external rhetorical situations modulate the loop from without---how the rhetor's \textit{logos} and the ambient \textit{Stimmung} of a civic assembly reach into the agent's chain and alter which action-type becomes likely---constitutes the central architectural question of the rhetoric chapter. Heidegger's observation that \textit{pathē} intrude upon \textit{logos} as articulated judgement (Struever, \textit{Heidegger and Rhetoric}, p.~109) and Burke's attention to how purposive orientation transforms indeterminate experience into motivated symbolic action (Burke, \textit{A Grammar of Motives}, pp.~311--312) will jointly furnish the external-modulation framework. Similarly, Gross's affordance-theoretic reconception of rhetoric as disclosing situations that are ``threatening and promising with respect to our human being-in-the-world'' will anchor the analysis of \textit{Mitsein}-mediated persuasion. \item \textbf{The \textit{De Anima} I.4, 408b5--7 cross-section warrant.} The direct application of the three-factor schema to cognitive and \textit{pathic} motions---rather than its merely analogical extension---bears not only on the present section but also on the cognitive-actuality and emotion-actuality analyses. Aristotle's insistence that it is not the soul but the human being who acts through the soul (Aristotle, \textit{De Anima}, 408b5--7) provides the textual warrant for treating all three domains---perceptual, cognitive, \textit{pathic}---as genuine instances of a single kinetic structure. The cross-section consistency of this warrant-deployment, while maintained throughout, requires explicit consolidation. \item \textbf{Integration with \textit{enargeia}.} The content-conditionality of doxastic ratification specifies that not every \textit{phantasma} produces higher-order \textit{pathē}; the \textit{phantasma} must engage evaluatively complex categories. The full mechanism by which \textit{enargeia}-induction moves an audience toward the concretion of higher-order \textit{pathē}, or conversely toward affectively inert routine through \textit{phantasma} de-evaluation, belongs to the rhetoric chapter's development. \end{enumerate}

Each of these questions converges on a single architectural concern: how external rhetorical situations reach into the agent's \textit{pathetic} loop, modulating which type of action becomes likely and hence reshaping the affective architecture of being-in-the-world.

%---END 09-DevNote---

\end{document}
